\documentclass[12pt]{article}
\usepackage[margin=1in]{geometry}
\usepackage{amsmath,booktabs,graphicx,threeparttable}

\begin{document}

\providecommand{\StructuralRobustnessPath}{appendix/structural-robustness}

\subsection{Sensitivity to Signal-Specific Social-Image Weights}
\label{app:signal-specific-lambda}

The baseline model restricts the social-image weight $\lambda$ to be common
across experimental arms. We assess this restriction using nested
specifications with a common weight, two weights pooling Ink and Bracelet as
Any Signal and Control and Calendar as No Signal, and four arm-specific
weights. The flexible specifications parameterize log-weight contrasts around
a common geometric mean. This keeps the common level of social-image concerns
separate from treatment-arm heterogeneity in those concerns.

Tables~\ref{tab:lambda-nested-ates} and
\ref{tab:lambda-nested-multipliers} show that allowing signal-specific weights
has little effect on the model-implied treatment effects or the qualitative
multiplier pattern. In all three specifications, the multiplier is above one
for Control and Calendar and below one for Ink and Bracelet. The flexible
models modestly move all multipliers toward one.

\begin{table}[htbp]
  \centering
  \small
  \caption{Structural treatment effects under alternative social-image weights}
  \label{tab:lambda-nested-ates}
  \begin{tabular}{llccc}
\toprule
Model & Treatment & Combined & Close & Far \\
\midrule
Common $\lambda$ & Bracelet & 0.08 & 0.06 & 0.12 \\
 &  & (0.06, 0.11) & (0.03, 0.08) & (0.09, 0.15) \\
 & Calendar & 0.03 & 0.02 & 0.04 \\
 &  & (0.01, 0.06) & (-0.01, 0.05) & (0.02, 0.07) \\
 & Ink & -0.02 & -0.05 & 0.02 \\
 &  & (-0.04, 0.01) & (-0.08, -0.02) & (-0.01, 0.05) \\
\midrule
Signal-group $\lambda$ & Bracelet & 0.08 & 0.06 & 0.11 \\
 &  & (0.06, 0.11) & (0.03, 0.09) & (0.08, 0.15) \\
 & Calendar & 0.03 & 0.02 & 0.04 \\
 &  & (0.00, 0.06) & (-0.01, 0.05) & (0.01, 0.07) \\
 & Ink & -0.02 & -0.04 & 0.01 \\
 &  & (-0.04, 0.01) & (-0.07, -0.01) & (-0.01, 0.05) \\
\midrule
Arm-specific $\lambda$ & Bracelet & 0.08 & 0.06 & 0.11 \\
 &  & (0.06, 0.11) & (0.03, 0.09) & (0.08, 0.14) \\
 & Calendar & 0.03 & 0.02 & 0.04 \\
 &  & (0.00, 0.06) & (-0.01, 0.05) & (0.01, 0.07) \\
 & Ink & -0.02 & -0.04 & 0.01 \\
 &  & (-0.05, 0.01) & (-0.08, -0.01) & (-0.02, 0.04) \\
\bottomrule
\end{tabular}

  \begin{minipage}{0.93\linewidth}
    \vspace{0.4em}\footnotesize
    \textit{Notes:} The common specification imposes one $\lambda$ across all
    arms. The signal-group specification estimates separate weights for Any
    Signal (Ink and Bracelet) and No Signal (Control and Calendar). The
    arm-specific specification estimates four weights. Parentheses contain
    95 percent posterior credible intervals.
  \end{minipage}
\end{table}

\begin{table}[htbp]
  \centering
  \small
  \caption{Social multipliers under alternative social-image weights}
  \label{tab:lambda-nested-multipliers}
  \begin{tabular}{llccc}
\toprule
Model & Treatment & 500m & 1500m & 2500m \\
\midrule
Common $\lambda$ & Control & 1.43 & 1.48 & 1.57 \\
 &  & (1.24, 1.73) & (1.24, 1.87) & (1.26, 2.03) \\
 & Ink & 0.92 & 0.86 & 0.81 \\
 &  & (0.80, 1.03) & (0.74, 0.97) & (0.70, 0.94) \\
 & Calendar & 1.26 & 1.22 & 1.21 \\
 &  & (1.11, 1.47) & (1.05, 1.46) & (1.01, 1.52) \\
 & Bracelet & 0.89 & 0.83 & 0.78 \\
 &  & (0.70, 1.04) & (0.69, 0.95) & (0.67, 0.90) \\
\midrule
Signal-group $\lambda$ & Control & 1.35 & 1.40 & 1.47 \\
 &  & (1.13, 1.71) & (1.15, 1.82) & (1.17, 1.98) \\
 & Ink & 0.93 & 0.88 & 0.84 \\
 &  & (0.79, 1.03) & (0.74, 0.98) & (0.70, 0.96) \\
 & Calendar & 1.21 & 1.18 & 1.18 \\
 &  & (1.07, 1.46) & (1.03, 1.45) & (1.01, 1.50) \\
 & Bracelet & 0.91 & 0.85 & 0.80 \\
 &  & (0.71, 1.02) & (0.71, 0.96) & (0.69, 0.92) \\
\midrule
Arm-specific $\lambda$ & Control & 1.35 & 1.40 & 1.47 \\
 &  & (1.14, 1.68) & (1.16, 1.78) & (1.17, 1.92) \\
 & Ink & 0.95 & 0.90 & 0.86 \\
 &  & (0.83, 1.02) & (0.77, 0.98) & (0.73, 0.96) \\
 & Calendar & 1.19 & 1.17 & 1.17 \\
 &  & (1.06, 1.42) & (1.04, 1.42) & (1.01, 1.45) \\
 & Bracelet & 0.91 & 0.86 & 0.81 \\
 &  & (0.72, 1.03) & (0.71, 0.96) & (0.68, 0.92) \\
\bottomrule
\end{tabular}

  \begin{minipage}{0.93\linewidth}
    \vspace{0.4em}\footnotesize
    \textit{Notes:} Entries are posterior medians of the social multiplier at
    the indicated distance. Parentheses contain 95 percent posterior credible
    intervals. Model definitions are given in the notes to
    Table~\ref{tab:lambda-nested-ates}.
  \end{minipage}
\end{table}

Because extra weights can be weakly identified, we vary the prior standard
deviation of pairwise log-weight ratios. Table~\ref{tab:lambda-prior} reports
the resulting contrast and multiplier estimates at 1.5 km. The posterior
contrasts remain near zero, and the multiplier conclusions are stable across
the prior grid. The contraction statistic is reported to distinguish a stable
posterior caused by data from one inherited mechanically from the prior.

\begin{table}[htbp]
  \centering
  \scriptsize
  \caption{Prior sensitivity of signal-specific social-image weights}
  \label{tab:lambda-prior}
  \resizebox{\linewidth}{!}{%
    \begin{tabular}{llrrrrrr}
\toprule
Model & Prior SD & $\log(\lambda_S/\lambda_N)$ & Contraction & Control & Ink & Calendar & Bracelet \\
\midrule
Grouped & 0.10 & 0.02 & 0.51 & 1.43 & 0.87 & 1.19 & 0.84 \\
Grouped & 0.25 & 0.04 & 0.32 & 1.40 & 0.88 & 1.18 & 0.85 \\
Grouped & 0.50 & 0.04 & 0.19 & 1.41 & 0.88 & 1.18 & 0.85 \\
Grouped & 1.00 & 0.04 & 0.11 & 1.39 & 0.88 & 1.17 & 0.85 \\
Arm-specific & 0.10 & 0.01 & 0.72 & 1.43 & 0.88 & 1.19 & 0.85 \\
Arm-specific & 0.25 & -0.01 & 0.57 & 1.40 & 0.90 & 1.17 & 0.86 \\
Arm-specific & 0.50 & -0.03 & 0.45 & 1.36 & 0.91 & 1.16 & 0.87 \\
Arm-specific & 1.00 & -0.05 & 0.30 & 1.33 & 0.92 & 1.14 & 0.88 \\
\bottomrule
\end{tabular}
}
  \begin{minipage}{0.96\linewidth}
    \vspace{0.4em}\footnotesize
    \textit{Notes:} Prior SD is the standard deviation assigned to pairwise
    log-$\lambda$ ratios. The contrast compares Any Signal with No Signal.
    Multiplier columns are evaluated at 1.5 km. Contraction reports the
    posterior standard deviation relative to its induced prior standard
    deviation; smaller values indicate greater learning.
  \end{minipage}
\end{table}

Figure~\ref{fig:lambda-prior-profile} complements the posterior comparison
with the regularized profile support for the Any Signal--No Signal contrast.
The profile retains the baseline priors on nuisance parameters and therefore
is not interpreted as a likelihood-ratio confidence interval.

\begin{figure}[htbp]
  \centering
  \includegraphics[width=.48\linewidth]{\StructuralRobustnessPath/figures/main-core-lambda-prior-sensitivity.pdf}
  \hfill
  \includegraphics[width=.48\linewidth]{\StructuralRobustnessPath/figures/main-core-lambda-profile.pdf}
  \caption{Prior sensitivity and regularized profile support for signal-specific image weights}
  \label{fig:lambda-prior-profile}
\end{figure}

Finally, a simulation exercise fits the common, grouped, and arm-specific
models to data generated under common, grouped, arm-heterogeneous, and
no-image designs. Table~\ref{tab:lambda-recovery} reports recovery of the image
weight and multiplier. The grouped and arm-specific models recover genuine
heterogeneity substantially better than the common-weight restriction, while
the more flexible specifications do not generate large multiplier errors when
the true weight is common.

\begin{table}[htbp]
  \centering
  \scriptsize
  \caption{Simulation recovery of signal-specific social-image weights}
  \label{tab:lambda-recovery}
  \begin{tabular}{llrrrr}
\toprule
DGP & Fitted model & $N$ & Bias($\lambda$) & RMSE($\lambda$) & RMSE(multiplier) \\
\midrule
arm-heterogeneous & arm & 50 & -0.099 & 0.230 & 0.119 \\
arm-heterogeneous & common & 50 & 0.474 & 0.533 & 0.181 \\
arm-heterogeneous & grouped & 50 & -0.151 & 0.292 & 0.122 \\
common & arm & 50 & -0.239 & 0.296 & 0.109 \\
common & common & 50 & -0.203 & 0.257 & 0.105 \\
common & grouped & 50 & -0.231 & 0.287 & 0.112 \\
grouped-025 & arm & 50 & -0.057 & 0.177 & 0.102 \\
grouped-025 & common & 50 & 1.107 & 1.132 & 0.293 \\
grouped-025 & grouped & 50 & -0.061 & 0.165 & 0.102 \\
grouped-050 & arm & 50 & 0.131 & 0.204 & 0.119 \\
grouped-050 & common & 50 & 2.297 & 2.318 & 1.132 \\
grouped-050 & grouped & 50 & 0.106 & 0.157 & 0.114 \\
no-image & arm & 50 & 0.072 & 0.119 & 0.049 \\
no-image & common & 50 & 0.039 & 0.076 & 0.031 \\
no-image & grouped & 50 & 0.079 & 0.128 & 0.053 \\
\bottomrule
\end{tabular}

  \begin{minipage}{0.93\linewidth}
    \vspace{0.4em}\footnotesize
    \textit{Notes:} Each row summarizes 50 simulated datasets. Bias and RMSE
    for $\lambda$ concern the relevant signal-specific contrast; the final
    column reports RMSE of the implied social multiplier.
  \end{minipage}
\end{table}

\begin{figure}[htbp]
  \centering
  \includegraphics[width=.60\linewidth]{\StructuralRobustnessPath/figures/main-core-lambda-recovery-contrast.pdf}
  \caption{Simulation recovery of the signal-specific social-image contrast}
  \label{fig:lambda-recovery}
\end{figure}

\end{document}
